\section{Discussion}
\label{sec:discussion}

\paragraph{What each mechanism actually buys.}
Legend shrinking and confidence rejection both trade coverage for macro quality, but they charge
different customers. Shrinking the legend refuses an entire crop type everywhere: a farmer growing
it receives nothing, in every parcel, in every region. Abstention refuses individual parcels
scattered across all crop types, so no grower is systematically excluded, which is the selective
classification literature's own concern about who pays for the abstention
\cite{jones2020selectivedisparities}. Our measurements say that abstention is also the better buy
at equal coverage. The two arguments point the same way, and we know of no operational setting in
which the reverse is preferable except one: when the downstream consumer requires a fixed,
published catalog and cannot act on a partial map.

\paragraph{What should be reported.}
A macro-averaged figure computed over a pruned catalog is uninterpretable on its own, and the fix
is neither expensive nor novel. Report the number of classes the average is taken over, the
coverage the product delivers, and the control of Sect.~\ref{sec:control}. Three numbers, one of
which costs a single line of code. Without them, a paper that raises macro F1 by a quarter of a
point after retiring three rare classes is reporting arithmetic, and a reader cannot tell it from a
paper that improved a model.

\paragraph{On the criterion practitioners use.}
The finding that retirement by support is both the weakest criterion and non-monotone deserves
emphasis, because support is the criterion that gets used. It is chosen for a defensible reason ---
a class with few samples is a class whose estimate is unstable --- but instability of the estimate
and poor separability are different properties, and support only tracks the first. A class with
\num{1718} parcels that the model separates cleanly contributes more to a delivered map than two
abundant classes the model cannot tell apart. Choosing by measured per-class quality, on blocks
disjoint from those that will be evaluated, costs nothing extra and is strictly better in every
comparison we ran.

\paragraph{On measuring meta-learners.}
The three figures for the same stacking model --- \num{0.7470}, \num{0.6789} and \num{0.5360} ---
are not three plausible estimates to choose among. The first is contaminated, and Wolpert's
original formulation of stacking already required the leakage-free partitions it violates
\cite{wolpert1992stacking}. The third is leak-free but penalizes the model for an aggregation
artifact: averaging per-block macro scores in an imbalanced setting scores blocks over classes
their parcels do not contain. Pooling the out-of-fold posteriors and scoring once, which gives
\num{0.6789}, is the estimator that answers the question actually being asked. The gap between the
first and the last is wide enough to reverse a conclusion about whether the ensemble contributes
anything, which is why the regime has to be stated whenever such a number is printed.
